\begin{lemma}[The zeroth Hermite mode is the constant function $1$]
Fix $n \in \mathbb{N}$ and write $h_0$ for the Hermite function of index
$\alpha = 0$, regarded through \noderef{HermiteL2} as an element of
$L^2(\gamma_n)$, with $\gamma_n$ the Gaussian measure of \noderef{GaussMeasure}.
Then all of the following hold.
\begin{enumerate}
  \item For every function $f : \mathbb{R}^n \to \mathbb{R}$,
    \[
      \int_{\mathbb{R}^n} f(x)\, h_0(x) \, d\gamma_n(x)
        = \int_{\mathbb{R}^n} f(x) \, d\gamma_n(x) .
    \]
  \item For every $v \in L^2(\gamma_n)$,
    $\displaystyle \langle v, h_0\rangle
       = \int_{\mathbb{R}^n} v \, d\gamma_n$.
  \item For every $v \in L^2(\gamma_n)$,
    $\displaystyle \langle h_0, v\rangle
       = \int_{\mathbb{R}^n} v \, d\gamma_n$.
  \item $\langle h_0, h_0\rangle = 1$.
  \item $\|h_0\|_{L^2(\gamma_n)} = 1$.
\end{enumerate}
Here $\langle\cdot,\cdot\rangle$ is the real inner product of $L^2(\gamma_n)$.
In short: $h_0$ is the constant function $1$, so integrating against it is
integrating against $\gamma_n$, and it is a unit vector.

Two points about the form of the statement.  First, $h_0$ is an element of
$L^2(\gamma_n)$, that is, an equivalence class, so the identity $h_0 = 1$ holds
only $\gamma_n$-almost everywhere; every clause above is therefore obtained from
a congruence of integrands almost everywhere, and never from a pointwise
substitution.  Clause (1) carries no integrability hypothesis on $f$ for exactly
this reason: it is the equality of the integrals of two functions that agree
$\gamma_n$-almost everywhere, which holds whatever the integrability status of
$f$.  Second, (2) and (3) differ only by the symmetry of the real inner product
and are both recorded because consumers meet the pairing in both argument orders.
\end{lemma}
\begin{proof}
By \noderef{HermiteL2ZeroAEOne},
\begin{equation*}
  h_0(x) = 1 \qquad \text{for } \gamma_n\text{-almost every } x \in \mathbb{R}^n .
  \tag{$\ast$}
\end{equation*}
This is the only input about $h_0$ used below; everything else is formal.

\emph{Clause (1).}  Let $f : \mathbb{R}^n \to \mathbb{R}$.  By $(\ast)$ the two
integrands agree almost everywhere: for $\gamma_n$-almost every $x$ we have
$f(x)\,h_0(x) = f(x)\cdot 1 = f(x)$.  Integrals of functions that agree
$\gamma_n$-almost everywhere are equal, so
$\int f\,h_0 \, d\gamma_n = \int f \, d\gamma_n$.  Note that this is a congruence
of integrands and not a rewriting of $h_0$ by $1$ at a point, which $(\ast)$ does
not license.

\emph{Clause (2).}  Let $v \in L^2(\gamma_n)$.  By \noderef{L2InnerGaussIntegral}
the inner product of two elements of $L^2(\gamma_n)$ is the $\gamma_n$-integral of
their product, so
\[
  \langle v, h_0\rangle
    = \int_{\mathbb{R}^n} v(x)\, h_0(x) \, d\gamma_n(x)
    = \int_{\mathbb{R}^n} v \, d\gamma_n ,
\]
the second equality being clause (1) applied to $f = v$.

\emph{Clause (3).}  The inner product of the real Hilbert space $L^2(\gamma_n)$ is
symmetric, so $\langle h_0, v\rangle = \langle v, h_0\rangle$, and clause (2)
finishes.

\emph{Clause (4).}  Apply clause (2) with $v = h_0$, which is legitimate since
$h_0 \in L^2(\gamma_n)$:
\[
  \langle h_0, h_0\rangle = \int_{\mathbb{R}^n} h_0 \, d\gamma_n .
\]
By $(\ast)$ again the integrand agrees $\gamma_n$-almost everywhere with the
constant function $1$, so this integral equals
$\int_{\mathbb{R}^n} 1 \, d\gamma_n = \gamma_n(\mathbb{R}^n)$, which is $1$
because $\gamma_n$ is a probability measure by \noderef{GaussIsProb}.  Hence
$\langle h_0, h_0\rangle = 1$.

\emph{Clause (5).}  In a real inner product space $\langle w, w\rangle = \|w\|^2$,
so clause (4) gives $\|h_0\|^2 = 1$.  Since $\|h_0\| \geq 0$ and the only
nonnegative real whose square is $1$ is $1$ itself, $\|h_0\| = 1$.
\end{proof}
